\begin{exercise}[Equilibres triphasés de l'ammoniac]
%Psi, Dunod, pg 67
La pression de vapeur saturante, notée $P_v^\star (T)$, de l'ammoniac liquide est donnée en fonction de la température par la relation ($P$ en pascal, $T$ en kelvin) :
\begin{align}
\ln P_v^\star (T) = 19.49 - \frac{3063}{T}
\end{align}
et la pression de sublimation, notée $P_s^\star (T)$, de l'ammoniac solide par :
\begin{align}
\ln P_s^\star (T) = 23.03 - \frac{3754}{T}
\end{align}
\begin{questions}
\question En déduire la température du point triple de l'ammoniac.
\question Calculer les enthalpies standard de vaporisation, sublimation et fusion au point triple.
\end{questions}
\paragraph*{Données} : 
Constante des gaz parfaits : $R = 8.314 ~ \si{\joule\per\mole\per\kelvin}$

\end{exercise}